\begin{abstract}

面向算力增长放缓与高质量语料稀缺的双重约束，大语言模型训练需要同时协调参数规模、训练数据量、数据质量、领域配比和长上下文开销。本文建立“数据属性建模——广义标度律——算力约束优化——能力前沿预测”的分层框架，对不同来源数据采用拟合、验证、压力测试和外推分层处理，并仅在通过物理与数值验收后向后续问题传递冻结接口。

针对问题一，将 $22$ 项异构质量信号同向化后划分为五个语义块，利用稳定 CRITIC 权重与冲突自适应 Huber 聚合得到样本和领域质量；在单纯形上建立二阶 Scheff\'e--Ridge 多输出响应面，并在经验支持域内进行多起点近优配比搜索。质量聚合方案扰动下七域排序的最低 Kendall 系数为 $0.8095$；二阶 Ridge 内部 NRMSE 为 $0.3959$，在 $1\mathrm{M}$、$60\mathrm{M}$ 和 $1\mathrm{B}$ 实测检验上的 Spearman 系数分别为 $0.8582$、$0.8698$ 和 $0.8503$。近优解对正则强度较稳健，但对支持半径和领域份额上界敏感，因而将其报告为近优配比族而非唯一全局最优解。

针对问题二，先比较参数--数据可分幂律、总算力幂律和交互幂律，再用分块留出验证选择质量修正形式，建立固定参考配比的广义标度律。参数--数据可分模型的留一规模 RMSE 为 $0.000147$，GQ2 质量修正模型在 $69$ 个留出块上的等权 RMSE 为 $0.072$。在中位参考点上，质量提高 $0.10$ 的 Loss 改善约等价于参数量增加 $37.30\%$ 或训练数据量增加 $56.95\%$。由于 A、B 体系缺少配比绝对效应的共同校准，非零配比强度在部分情景中还会跌破不可约损失，故联合配比接口不进入主决策。

针对问题三，固定参考领域配比，将基础训练、质量提升和长上下文注意力开销纳入同一 FLOPs 预算，联合优化 $N$、$D$ 和 $Q_A$。利用 Loss 对 $D$ 严格递减的结构解析消去 $D$，结合粗网格、多起点局部精化、三维独立复核与 KKT 边际条件识别资源结构转移。注意力开销与基础训练开销相当的上下文临界长度为 $30000$ Token。在 $4096$ 上下文、指数质量成本和 $10^{22}$ FLOPs 下，最优配置为 $N^*=5.44886\mathrm{B}$、$D^*=244.275\mathrm{B}$、$Q_A^*=0.999999$，预测 Loss 为 $2.15491$。连续预算路径呈现“质量未启动——质量与规模竞争——质量饱和后规模重新主导”的阶段变化，高预算结果则需保留外推标记。

针对问题四，以六项 Benchmark 重建综合能力，通过单调 Loss--Benchmark 桥接将问题三的 $45$ 个冻结情景映射到能力尺度，再以条件分位动力学模型和 Shapley 路径平均分解历史前沿。$2024$ 年 $6$ 月至 $2025$ 年 $1$ 月的能力增量为 $4.6182$ 分，其中规模扩张相关贡献为 $4.6183$ 分，控制规模后的非规模残余区间跨零，尚无稳定方向。若近期算力对数增速分别保留 $100\%$、$50\%$ 和 $25\%$，则 $24$ 个月后 $0.90$ 能力前沿中心值约为 $77.15$、$65.74$ 和 $59.16$。上述结果均是给定数据截止日、支持域和成本情景下的条件结论。

\noindent \textbf{关键词：}大语言模型\quad 数据质量\quad 广义标度律\quad 算力约束\quad 资源配置\quad 能力前沿

\end{abstract}
